\section{Introduction}\label{sec:intro}

Inverted pendulums are among the most effective passive low-frequency mechanical filters. In the seismic attenuation chains of interferometric gravitational-wave detectors, rigid legs hold the payload in the unstable upright configuration while an elastic element is tuned to almost cancel the gravitational antispring, so that the resonance frequency falls far below the seismic band and ground motion is strongly suppressed~\cite{Losurdo1999,Takamori2007}. The spring-coupled inverted pendulum (PIAR, the acronym of its Spanish name, \textit{p\'endulo invertido acoplado a un resorte}) captures this principle in its simplest form: a rigid rod pivoted at its lower end carries a mass and is held upright by a horizontal spring whose far end rides on a collar that slides along a vertical guide. Its conservative dynamics is controlled by a single dimensionless number, $\mu=kL/(Mg)-1$, which measures how much the elastic restoring torque exceeds the gravitational toppling torque. As $\mu\to0^{+}$ the small-oscillation frequency vanishes, which is the property sought for isolation, but the energy barrier that protects the upright position vanishes even faster, which is its price. Once damping and periodic forcing are added, the same model displays, in its weak-spring regime ($\mu<0$), the period-doubling route to chaos familiar from the driven pendulum~\cite{BakerGollub1996,Strogatz2015,GuckenheimerHolmes1983}. Inverted pendulums also remain paradigmatic models of stability and of its dynamic control~\cite{Butikov2001}. Practical relevance, analytical tractability, and dynamical richness thus meet in a single degree of freedom.

Practical relevance, however, comes with a practical difficulty. The state of an isolator is rarely measured completely: angle sensors, load cells, and position encoders drop out, saturate, or are simply absent, and the unmeasured variables must be inferred from the measured ones. Statistics treats this situation as a missing-data problem~\cite{LittleRubin2019}, and deep imputation methods for time series, from bidirectional recurrent networks~\cite{Cao2018} to Gaussian-process variational autoencoders~\cite{Fortuin2020} and score-based diffusion models~\cite{Tashiro2021}, have become remarkably accurate on generic benchmarks. For deterministic, and in particular chaotic, dynamics a different and more physical route exists. Delay coordinates of a single observable reconstruct the attractor of the system~\cite{Packard1980,Takens1981}, generically and up to a smooth change of coordinates~\cite{Sauer1991}. Nonlinear time-series analysis and prediction~\cite{Abarbanel1993,KantzSchreiber2004,Casdagli1989,Pathak2018}, the detection of causality by cross mapping between reconstructed manifolds~\cite{Sugihara2012}, and the inference of unmeasured variables with reservoir computers~\cite{Lu2017} all rest on this fact, which has become a main point of contact between machine learning and complex systems~\cite{Tang2020}.

The full--partial reconstruction mapping (FPRM), recently introduced by Wu et al.~\cite{Wu2026FPRM}, turns this geometric fact into an imputation method. The delay reconstruction of a completely observed time series and that of a partially observed one are both faithful images of the same attractor; the map between them can therefore be learned at the times at which both series are available and evaluated wherever the partial series is missing. Implemented with Gaussian-process regression (GPR)~\cite{Rasmussen2006}, the FPRM recovers complex dynamics even when most of the data are missing, and it rests on a transparent justification rooted in dynamical-systems theory. Two questions left open by the original formulation become pressing for mechanical systems. The first concerns identifiability: Takens' theorem holds for generic observables, but physical symmetries can make a measured quantity nongeneric~\cite{LetellierAguirre2002,Letellier2005}. The missing value is then no longer a function of the delay coordinates, its conditional distribution becomes multimodal, and a Gaussian regressor returns the average of the branches, a value that the system never visits. The second question concerns physics: the equations of motion of a mechanical system are usually known, at least in form, and it is natural to ask whether imposing them improves the reconstruction.

Conditional generative models address the first question by construction. Normalizing flows~\cite{Papamakarios2021,Durkan2019}, variational autoencoders~\cite{Kingma2014,Sohn2015}, and denoising diffusion models~\cite{Ho2020,NicholDhariwal2021} learn full conditional distributions and can represent several branches at once. Physics-informed learning~\cite{Raissi2019,Karniadakis2021} addresses the second question by adding the residuals of known equations to the training loss. Structure-preserving architectures, such as Hamiltonian and Lagrangian neural networks~\cite{Greydanus2019,Cranmer2020} and symplectic networks~\cite{Jin2020}, instead build conservation laws into the model itself. Soft physical constraints are known, however, to reshape the loss landscape in ways that can impair training~\cite{Krishnapriyan2021,WangTengPerdikaris2021}, and whether they help a generative imputer is an open empirical question.

The PIAR is a well-suited test bed for both questions. Its invariants (equilibria, barrier, separatrix energies, and normal form) are available in closed form, and its periods as quadratures, so that numerical and learned results can be checked against exact mechanics. Its two natural measurements, the spring force and the height of the collar, are nontrivial functions of the angle that cannot be converted into each other by differentiation, and the reflection $\theta\to-\theta$ makes the first odd and the second even, which creates a controlled, physically motivated identifiability problem. Its bifurcation parameter provides a single axis along which the dynamics changes qualitatively, and, under forcing, the system develops a strange attractor of low and measurable dimension.

This article develops the PIAR into a benchmark for dynamics-informed reconstruction and uses it to extend the FPRM in three respects: a symmetry analysis, verified by a model-free cross-mapping test, that establishes in advance which reconstruction tasks are identifiable; a generative engine, a conditional normalizing flow, that represents the multimodal distributions that arise when identifiability fails; and a controlled assessment of when equation-of-motion residuals improve the reconstruction and when they induce mode collapse. Section~\ref{sec:model} presents the analytical mechanics in a self-contained form, and Sec.~\ref{sec:cases} uses high-accuracy integration and exact periods to characterize five canonical regimes of the conservative dynamics, which span the three stability classes of the upright position. Section~\ref{sec:chaos} follows the damped--driven dynamics through a Feigenbaum cascade to a strange attractor, uses Melnikov's method to locate the splitting of the separatrix from which chaos grows, and relates the dimension of the attractor to the embedding dimension required by Takens' theorem. Section~\ref{sec:fprm} establishes, by a model-free cross-mapping analysis, which reconstruction tasks are well posed and how the phase of the forcing restores the identifiability broken by the reflection symmetry; on this basis, the Gaussian-process engine of the FPRM is replaced by a conditional neural spline flow, optionally regularized by equation-of-motion residuals. Section~\ref{sec:results} benchmarks the resulting generative FPRM against the original method and against variational and diffusion alternatives. Section~\ref{sec:lab} describes the open, interactive laboratory that regenerates every figure and number of this article from fixed random seeds, Sec.~\ref{sec:discussion} discusses the implications and limitations of the results, and Sec.~\ref{sec:conclusions} summarizes the conclusions.
